\section*{第 \ref{ch:g10:layered-ocean} 章——分层的海洋}

\begin{solution}{exo:g10:layered-ocean:1}
\omterm{def:g10:layered-ocean:mixed-layer}{混合层}是海洋表层的一层，其中温度和
盐度几乎不随深度变化；\omterm{def:g10:layered-ocean:thermocline}{温跃层}是位于其下方的一层，
其中温度随深度迅速下降。风和波浪搅动
\omterm{def:g10:layered-ocean:mixed-layer}{混合层}（冬季时海表降温也有助于此）。
\end{solution}

\begin{solution}{exo:g10:layered-ocean:2}
\omterm{def:g10:layered-ocean:upwelling}{上升流}是来自表层以下的冷水上涌，取代
被风吹离的表层水。上升的水体将矿物营养物质
带入有光层，浮游植物利用它们生长；
浮游植物为食物链的其他环节提供食物。
\end{solution}

\begin{solution}{exo:g10:layered-ocean:3}
热带太平洋：表层约 \qty{29}{\degreeCelsius}，
深 \qty{1000}{m} 处约 \qty{4.5}{\degreeCelsius}。北冰洋：
表层约 \qty{-1.4}{\degreeCelsius}，
深 \qty{1000}{m} 处接近 \qty{0}{\degreeCelsius}。
\end{solution}

\begin{solution}{exo:g10:layered-ocean:4}
在 \qty{0}{\degreeCelsius} 时（\qty{1028.1}{kg/m^3}，对比
\qty{1023.3}{kg/m^3}，\qty{25}{\degreeCelsius} 时）。海水：在 \qty{4}{\degreeCelsius} 时，纯水
的密度为 \qty{1000.0}{kg/m^3}，盐度为 35 的海水约为
\qty{1027.8}{kg/m^3}（由图中读出）。
\end{solution}

\begin{solution}{exo:g10:layered-ocean:5}
浮标在 \qty{1000}{m} 处漂流，下沉至 \qty{2000}{m}，然后上升至
表层，途中测量温度、盐度（由传导率算得）和压强；
到达表层后，通过卫星发送其剖面数据，然后再次下沉。
浮标随洋流在全球海洋中漂移，可连续工作数年，
各季节均可工作，且远离船舶航线。
\end{solution}

\begin{solution}{exo:g10:layered-ocean:6}
$361.84\times10^6 \times 3.682 \approx \qty{1.33e9}{km^3}$，约
\num{1.33} 十亿立方千米。$100/3682 \approx 0.027$：不足
平均深度的 \qty{3}{\percent}。
\end{solution}

\begin{solution}{exo:g10:layered-ocean:7}
$22.7 - 0.5 = \qty{22.2}{\degreeCelsius}$，该值在 \qty{15}{m} 处达到，\omterm{def:g10:layered-ocean:mixed-layer}{混合
层}深约 \qty{15}{m}。\qty{20}{m} 至 \qty{50}{m} 之间：
$(21.8 - 17.3)/30 = \qty{0.15}{\degreeCelsius}$ 每米。
\end{solution}

\begin{solution}{exo:g10:layered-ocean:8}
太阳只加热表层，季节性的加热和冷却
仅通过\omterm{def:g10:layered-ocean:mixed-layer}{混合层}扩散；\omterm{def:g10:layered-ocean:thermocline}{温跃层}以下的海水
与之隔绝。冬季表层降温、密度增大，下沉到
下层海水中，且风暴搅动海水，混合可深入。夏季
阳光加热轻质的表层，该表层漂浮在较冷的水体之上；
较弱的风无法将其压入深层，\omterm{def:g10:layered-ocean:mixed-layer}{混合层}因此保持较薄。
\end{solution}

\begin{solution}{exo:g10:layered-ocean:9}
$0.78/0.18 \approx 4.3$ 度，在 \qty{10}{\degreeCelsius} 附近；
$0.78/0.06 = 13$ 度，在 \qty{0}{\degreeCelsius} 附近。在寒冷的极地海域，
盐度的微小变化对密度的影响远大于温度变化，
在那里，\omterm{def:g10:layered-ocean:stratification}{层结}由盐度控制。
\end{solution}

\begin{solution}{exo:g10:layered-ocean:10}
温度相同：较深的海水更咸，因此密度更大
（盐度为 33.3，对比 28.1）。两者之间的一层，其中盐度
随深度迅速变化，称为\omterm{def:g10:layered-ocean:thermocline}{盐跃层}；此处它同时也是\omterm{def:g10:layered-ocean:thermocline}{密度跃层}。
\end{solution}

\begin{solution}{exo:g10:layered-ocean:11}
水团被压下后，其温度和盐度保持不变，因而密度也不变。
周围是密度更大的海水，它比其所排开的海水更轻：
浮力将其推回原来的深度。若将水团向上推，
则其密度将大于新处的周围海水，并沉回原位。这些分层是稳定的。
\end{solution}

\begin{solution}{exo:g10:layered-ocean:12}
质量：$0.001 \times 1028.1 = \qty{1.0281}{kg}$。新的体积：$1.0281/1023.3
\approx \qty{0.0010047}{m^3}$，约 \qty{1004.7}{mL}：体积增大了
约 \qty{4.7}{mL}，即 \qty{0.47}{\percent}。厚 \qty{700}{m} 的一层
将增厚 $700 \times 0.0047 \approx \qty{3.3}{m}$。
\end{solution}

\begin{solution}{exo:g10:layered-ocean:13}
这是二月的北大西洋。该处顶部 \qty{150}{m} 的水层几乎具有
相同的温度，因而密度也几乎相同：搅动它所需的能量很少。
在热带太平洋中，密度增加了 $1026.0 - 1021.6 =
\qty{4.4}{kg/m^3}$，出现在 \qty{50}{m} 至 \qty{200}{m} 之间：\omterm{def:g10:layered-ocean:thermocline}{密度跃层}很强，
风暴需要将轻水压入
密度大得多的海水中。
\end{solution}

\begin{solution}{exo:g10:layered-ocean:14}
远离海岸处叶绿素含量高：太平洋中沿赤道的带状区域
（大西洋中亦然），信风驱动\omterm{def:g10:layered-ocean:upwelling}{上升流}；以及
高纬度海域（如北大西洋），冬季\omterm{def:g10:layered-ocean:mixed-layer}{混合层}
深入并将营养物质带向表层。叶绿素含量极低的区域是亚热带环流
中心，\omterm{def:g10:layered-ocean:stratification}{层结}很强，没有任何过程将深层水带
上来，有光层的营养物质已耗尽。
\end{solution}

\begin{solution}{exo:g10:layered-ocean:15}
下沉：$1000/(6 \times 3600) \approx \qty{0.046}{m/s}$，约
\qty{4.6}{cm/s}。上升：$2000/(10 \times 3600) \approx \qty{0.056}{m/s}$，
约 \qty{5.6}{cm/s}。$36\,000/2 = 18\,000$ 次测量，平均
间隔 $2000/18\,000 \approx \qty{0.11}{m}$。
\end{solution}

\begin{solution}{pb:g10:layered-ocean:1}
\textbf{1.} 环流区高出 $24.0 - 19.0 = \qty{5.0}{\degreeCelsius}$。

\textbf{2.} 环流（\qty{20.5}{\degree S}，对照 \qty{13.5}{\degree S}）。
离赤道越远应当越冷：仅凭纬度
无法解释为何它会高出 \qty{5}{\degreeCelsius}。

\textbf{3.} $24.0 - 0.5 = \qty{23.5}{\degreeCelsius}$，该温度出现在
\qty{50}{m}（23.6）与 \qty{100}{m}（22.6）之间，约在 $50 + 50 \times
0.1/1.0 = \qty{55}{m}$ 处。

\textbf{4.} 否：温度从表层起持续下降，
在前 \qty{25}{m} 内下降 \qty{1.6}{\degreeCelsius}，前 \qty{100}{m} 内下降
\qty{5.3}{\degreeCelsius}；\omterm{def:g10:layered-ocean:thermocline}{温跃层}从表层即开始。

\textbf{5.} 在 \qty{200}{m} 与 \qty{300}{m} 之间，$20.3 - 14.4 =
\qty{5.9}{\degreeCelsius}$，每百米。

\textbf{6.} 秘鲁近岸在 \qty{50}{m} 与 \qty{100}{m} 之间：约
\qty{70}{m}（插值得出；完整剖面给出 \qty{65}{m}）。在
环流中，\qty{200}{m} 与 \qty{300}{m} 之间：约 $200 + 100 \times 5.3/5.9
\approx \qty{290}{m}$。

\textbf{7.} 是，两个水柱的各层均如此：密度较小的水位于
密度较大的水之上，两个水柱都是稳定的。

\textbf{8.} 盐度较高的水应当密度较大；然而环流表层水的
密度反而较小：\qty{1024.71}{kg/m^3}，另一处为 \qty{1025.08}{kg/m^3}。

\textbf{9.} 环流中的水盐度约高 1.3 个单位，使密度增加约
$1.3 \times 0.8 \approx \qty{1.0}{kg/m^3}$；温度则高出 \qty{5}{\degreeCelsius}，
使密度减小约 $5 \times 0.27 \approx \qty{1.4}{kg/m^3}$，在这些
温度下，温度起主导作用。

\textbf{10.} $1026.44 - 1025.08 = \qty{1.36}{kg/m^3}$。

\textbf{11.} $1025.30 - 1024.71 = \qty{0.59}{kg/m^3}$。

\textbf{12.} 从表层到 \qty{500}{m}，密度共增加
$1026.92 - 1024.71 = \qty{2.21}{kg/m^3}$，其中仅 \qty{0.59}{kg/m^3}
位于 \qty{200}{m} 以上：大部分，即 \qty{1.62}{kg/m^3}，位于
\qty{200}{m} 以下，即位于 \qty{200}{m} 至 \qty{500}{m} 之间的\omterm{def:g10:layered-ocean:thermocline}{温跃层}中。

\textbf{13.} 秘鲁近岸，密度增加 $1026.97 - 1025.08 =
\qty{1.89}{kg/m^3}$，从表层到 \qty{500}{m}，其中
\qty{1.36}{kg/m^3} 位于顶部 \qty{200}{m} 内：大部分靠近表层。

\textbf{14.} 在 \qty{500}{m} 处，两地几乎相同
（\qty{7.9}{\degreeCelsius} 与 \qty{7.1}{\degreeCelsius}，
\qty{1026.97}{kg/m^3} 与 \qty{1026.92}{kg/m^3}）：两处之下是同一团冷水；
差别在于它离表层的远近。

\textbf{15.} 沿海岸向赤道吹的风把表层水推向外海，
近岸的下层水上升以填补其位置。
\omterm{def:g10:layered-ocean:thermocline}{温跃层}中的冷水被带上来，因此\omterm{def:g10:layered-ocean:thermocline}{温跃层}
向海岸方向抬升，几乎升至表层。

\textbf{16.} 秘鲁近岸：温度为 \qty{13}{\degreeCelsius} 至
\qty{15}{\degreeCelsius} 的水来自\omterm{def:g10:layered-ocean:thermocline}{温跃层}以下，位于顶部
\qty{100}{m} 内，处于风力搅动与光照可及的范围内。在
环流中，它位于 \qty{250}{m} 以下，处于一层温暖、密度较小的水层之下，风
无法搅动到那么深。

\textbf{17.} \omterm{def:g10:layered-ocean:upwelling}{上升流}将停止：表层水会在
阳光照射下升温，顶部会形成温暖、密度较小的水层，
而向有光层输送的营养物质将减少。

\textbf{18.} 秘鲁海岸呈浅蓝至黄色（叶绿素含量高）；
南太平洋环流中心位于
图中颜色最深的蓝色区域之一（叶绿素极少）。

\textbf{19.} 浮游植物既需要营养物质，也需要光。在环流中，
强\omterm{def:g10:layered-ocean:stratification}{层结}使富含营养物质的深层水远在有光层之下，
表层的营养物质早已耗尽。

\textbf{20.} $1.36/0.59 \approx 2.3$：\textbf{秘鲁近岸顶部
\qty{200}{m} 内的密度增幅约为环流中的 2.3 倍。}
\omterm{def:g10:layered-ocean:thermocline}{密度跃层}及其下方富含营养物质的冷水，已被抬升
至秘鲁近岸的顶部 \qty{200}{m} 以内，而在环流中，它位于更深处，
在 \qty{200}{m} 以下。
\end{solution}
